% =====================================================================
\section{Background and Related Work}
% =====================================================================

\subsection{The Migration Problem}
The paper motivates the work with a running example: a proto message
describing a toy in an online store's inventory, where the
\texttt{toy\_id} field is declared as a 32-bit integer. As the number
of toys grows, the store risks exhausting the roughly 2.1 billion
values representable by a signed 32-bit integer, so the field is
widened to a 64-bit integer. This one-line schema edit ripples
outward, though: every place that reads or writes the field, every
bit-manipulation routine that assumed a 32-bit width, and every test
that constructs identifiers with small, in-range sample values needs
to be inspected and, in many cases, changed. Some references are
\emph{direct} (a method call that reads the field), while others are
\emph{indirect}, reached only through data flow, for instance a local
variable later passed into the field setter. The authors note that an
earlier, fully manual version of this same kind of migration took
roughly two years to complete for a single identifier, largely
because regular expressions used to find call sites were
simultaneously too broad (matching unrelated fields with similar
names) and too narrow (missing indirect references and magic numbers
tied to the old bit width).



\subsection{Code Translation and Transformation}
The paper situates its contribution within two long-standing lines of
programming-languages research. \emph{Code translation} concerns
converting code from one language to another, historically via
hand-written parsers and abstract-syntax-tree rewriting rules
\cite{cordy2004txl}, and more recently via learned translation models.
\emph{Code transformation} concerns rewriting code \emph{within} the
same language, traditionally using rule-based languages such as TXL
\cite{cordy2004txl} and Stratego/XT \cite{bravenboer2008stratego}, or
by inferring transformation-by-example rules from a small number of
sample edits. Recent systems such as CodePlan use LLMs with
program-analysis-guided planning to make repository-scale, multi-file
changes \cite{bairi2024codeplan}, and other work automatically mines
reusable transformation rules for library APIs from pull requests
\cite{ramos2023melt}. Closest in spirit are two smaller prior studies:
one evaluating GPT-4 on migrating a handful of methods off an old
version of the SQLAlchemy library \cite{almeida2024library}, and
another studying how developers collaborate with an AI
coding-migration assistant moving a Java 7 codebase to Java 18
\cite{tehrani2024evaluating}. The paper under reproduction differs
from both mainly in scale: it covers 39 distinct migrations, hundreds
of code changes, and tens of thousands of edits, carried out
continuously over a full year inside a production monorepo rather
than as a bounded, small-scale experiment.

\subsection{Automated Refactoring and Rule Mining}
Beyond translation and whole-file transformation, a broader strand of
work targets narrower, well-defined refactoring and transformation
tasks using LLMs paired with static or dynamic program analysis,
mining reusable transformation rules from examples rather than
requiring a developer to hand-write them
\cite{ramos2023melt, bairi2024codeplan}. Compared to this line of
work, the identifier bit-width migration studied in the reproduced
paper is unusual in that a single, uniform intent, widening an
identifier, must be applied consistently across an enormous number of
call sites in several different languages, rather than as a one-off
refactoring of a single method or class.

\subsection{Companion Reports at Google}
The authors note that this paper is a detailed follow-up to two
shorter, earlier write-ups describing the same system: an industry
practice report and an internal engineering blog post, both of which
introduced the system at a higher level before this paper's
twelve-month, 39-identifier case study was available.

The system builds on several pieces of existing Google infrastructure:
a single, monolithic source-code repository shared across the company
\cite{potvin2016google}; Kythe, a code-indexing system used to resolve
symbol references across languages; Critique, Google's internal,
browser-based code review tool \cite{sadowski2018modern}; and an
internally fine-tuned variant of the Gemini family of models
\cite{team2023gemini}, trained using the DIDACT framework for
software-engineering tasks such as code-review comment resolution and
build-error repair.